\exercisesection{HILBERT-SAMUEL SERIES}

\begin{enumerate}
\def\labelenumi{\arabic{enumi})}
\tightlist
\item
  Let \(\Gamma\) be a commutative group and let H be a graded ring of type \(\Gamma\) (A, II, p.~164).
\end{enumerate}

For \(\gamma\in\Gamma\) , we denote by H \(_{\gamma}\) the homogeneous component of degree \(\gamma\) of H. We assume that H is generated by H \(_{0}\) and a finite family (x \(_{i}\) ) \(_{i\in I}\) of homogeneous elements of degree \(\neq0\) . We denote by \(\gamma_{i}\) the degree of x \(_{i}\) . For every commutative group K, we denote by K{[}{[}Γ{]}{]} the product group K \(^{\Gamma}\) and we denote by \(\sum_{\gamma\in\Gamma}m_{\gamma}\text{T}^{\gamma}\) the element (m \(_{\gamma}\) ) of K \(^{\Gamma}\) . We denote by K{[}Γ{]} the subgroup K \(^{(\Gamma)}\) of K{[}{[}Γ{]}. We endow K{[}{[}Γ{]}{]} with its natural structure as a module over the algebra Z{[}Γ{]} of the group Γ. Finally, let C be a hereditary set of classes of H₀-modules (A, VIII, § 10, no. 1, def. 1) and let K(C) be the corresponding Grothendieck group (loc. cit., no. 2). For every H₀-module E of type C, we denote {[}E{]} its class in K(C).

\begin{enumerate}
\def\labelenumi{\alph{enumi})}
\tightlist
\item
  Suppose \(\mathbf{H}_0\) is Noetherian, and let \(\mathbf{M}\) be a finitely generated graded H-module such that for every \(\gamma \in \Gamma\) , the \(\mathbf{H}_0\) -module \(\mathbf{M}_{\gamma}\) is of type C. Show that the element \(\prod_{i \in I} (1 - T^{\gamma_i}) \sum_{\gamma \in \Gamma} [\mathbf{M}_{\gamma}] T^{\gamma}\) of
\end{enumerate}

\(K(\mathbb{C})[[\Gamma]]\) belongs to \(K(\mathbb{C})[\Gamma]\) .

\begin{enumerate}
\def\labelenumi{\alph{enumi})}
\setcounter{enumi}{1}
\tightlist
\item
  Suppose now that \(\Gamma = \mathbf{Z}^s\) , \(\mathrm{H}_{\gamma} = 0\) and \(\mathbf{M}_{\gamma} = 0\) for \(\gamma \notin \mathbf{N}^s\) , and that the \(\gamma_i\) are basis vectors of \(\mathbf{Z}^s\) ; denoting by \(a_k\) the number of \(\gamma_i\) that are equal to the \(k\) th basis vector, prove that one has
\end{enumerate}

\[
\prod_ {k = 1} ^ {s} (1 - \mathrm{T} _ {k}) ^ {a _ {k}} \sum_ {\gamma \in \mathbf {Z} ^ {s}} [ \mathrm{M} _ {\gamma} ] \mathrm{T} ^ {\gamma} = \sum_ {\gamma} m _ {\gamma} \mathrm{T} ^ {\gamma}
\]

where \(\mathbf{T} = (\mathrm{T}_{1},\dots,\mathrm{T}_{s})\) , and \((m_{\gamma})\) is a finitely supported family. Deduce the existence of elements \(c_{\boldsymbol{b}}\in\mathrm{K}(\mathbb{C})\), for \(\boldsymbol{b}\in\mathbf{Z}^{s}\), all but finitely many of which are zero, such that, for \(\gamma\in\mathbf{N}^{s}\) sufficiently large, one has

\[
[ \mathbf {M} _ {\gamma} ] = \sum_ {\boldsymbol {b}} c _ {\boldsymbol {b}} \mathbf {B} _ {\boldsymbol {b}} (\gamma)
\]

where

\[
\mathrm{B} _ {b} (\gamma) = \prod_ {j = 1} ^ {s} \frac {(\gamma_ {j} + 1) \dots (\gamma_ {j} + b _ {j})}{b _ {j} !},
\]

the sum being taken over the elements b of \(Z^{s}\) satisfying \(b_{i} < a_{i}\) for \(1 \leqslant i \leqslant s\) .

\begin{enumerate}
\def\labelenumi{\arabic{enumi})}
\setcounter{enumi}{1}
\tightlist
\item
  Let \(f: \mathbf{Z} \to \mathbf{Z}\) be a map. Show that the following three properties are equivalent:
\end{enumerate}

\begin{enumerate}
\def\labelenumi{\alph{enumi})}
\item
  There exists \(n_0, n_1\) in \(\mathbf{Z}\) and \(\mathrm{P}\) in \(\mathbf{Q}[\mathrm{T}]\) such that \(f(n) = 0\) for \(n \leqslant n_1\) and \(f(n) = \mathrm{P}(n)\) for \(n \geqslant n_0\) .
\item
  There exists a finitely supported family \((a_i)_{i\in \mathbf{Z}}\) of elements of \(\mathbf{Z}\) such that \(f(n) = \sum_{i\in \mathbf{Z}}a_i\left[ \begin{array}{cc}n + i\\ i \end{array} \right]\) for every \(n\in \mathbf{Z}\) .
\item
  There exists \(r \in \mathbf{Z}\) and \(Q \in \mathbf{Z}[T, T^{-1}]\) such that \(\sum_{n \in \mathbf{Z}} f(n) T^n = (1 - T)^r Q(T)\) .
\end{enumerate}

Then we say that \(f\) is polynomial at infinity.

Generalize to the case of maps from Z into an arbitrary Z-module.

\begin{enumerate}
\def\labelenumi{\arabic{enumi})}
\setcounter{enumi}{2}
\tightlist
\item
  \begin{enumerate}
  \def\labelenumii{\alph{enumii})}
  \tightlist
  \item
    Give an example of a graded algebra H of type Z, over a field K, generated by a finite number of homogeneous elements of degrees \textgreater{} 0, and such that the map \(n \mapsto \dim_{\mathbf{K}} \mathrm{H}_{n}\) is not polynomial at infinity.
  \end{enumerate}
\end{enumerate}

\begin{enumerate}
\def\labelenumi{\alph{enumi})}
\setcounter{enumi}{1}
\tightlist
\item
  With the notations and hypotheses of th. 1 of § 4, no. 2, let D be the lcm of the \(d_{i}\). Show that H is a finitely generated module over the ring \(\mathrm{H}_{0}[(\mathrm{X}_{i}^{\mathrm{D}/d_{i}})_{i\in\mathbb{I}}]\) .
\end{enumerate}

Deduce that for every \(n_{0}\in\mathbf{N}\) the series \(\sum_{n\in\mathbf{Z}}\mathrm{long}_{\mathrm{H}_{0}}(\mathbf{M}_{n_{0}+n\mathbf{D}})\mathbf{T}^{n}\) is of the form \(\frac{\mathbf{P}_{n_{0}}(\mathbf{T})}{(1-\mathbf{T})^{c}}\) where \(c=\mathrm{Card}(\mathbf{I})\) and \(\mathbf{P}_{n_{0}}(\mathbf{T})\in\mathbf{Q}[\mathbf{T},\mathbf{T}^{-1}]\) . Deduce that the map \(n\mapsto\mathrm{long}_{\mathrm{H}_{0}}(\mathbf{M}_{n_{0}+n\mathbf{D}})\) is polynomial at infinity.

\begin{enumerate}
\def\labelenumi{\arabic{enumi})}
\setcounter{enumi}{3}
\tightlist
\item
  Let us place ourselves in the situation of Th. 1 of § 4, No. 2, and denote by C the set of isomorphism classes of finitely generated graded H-modules with components of finite length over \(\mathbf{H}_0\), and by \(\mathrm{K}(\mathbb{C})\) the corresponding Grothendieck group.
\end{enumerate}

\begin{enumerate}
\def\labelenumi{\alph{enumi})}
\tightlist
\item
  For every H-module M as above, show that there exists an \(H_{0}\)-graded module of finite length \(M_{0}\) and a surjective graded homomorphism \(p:H\otimes_{H_{0}}M_{0}\to M\). Deduce from this that every H-module M as above admits a graded resolution (A, X, p.~56) by H-modules:
\end{enumerate}

\[
\dots \rightarrow \mathrm{H} \otimes_ {\mathrm{H} _ {0}} \mathrm{M} _ {1} \rightarrow \mathrm{H} \otimes_ {\mathrm{H} _ {0}} \mathrm{M} _ {0} \rightarrow \mathrm{M} \rightarrow 0
\]

where, for all \(i\), \(M_{i}\) is an \(H_{0}\)-graded module of finite length.

\begin{enumerate}
\def\labelenumi{\alph{enumi})}
\setcounter{enumi}{1}
\item
  Using a Koszul resolution of the H-module \(\mathbf{H}_0\) (A, X, p.~152), show that for every graded H-module M, \(\mathrm{Tor}_j^{\mathrm{H}}(\mathrm{H}_0,\mathrm{M})\) is canonically endowed with a grading, that \(\mathrm{Tor}_j^{\mathrm{H}}(\mathrm{H}_0,\mathrm{M}) = 0\) for \(j > \mathrm{Card}(\mathrm{I})\), and that \(\mathrm{Tor}_j^{\mathrm{H}}(\mathrm{H}_0,\mathrm{P}\otimes_{\mathrm{H}_0}\mathrm{H})\) is an \(\mathbf{H}_0\)-module of finite length for every \(j\) when P is an \(\mathbf{H}_0\)-module of finite length.
\item
  Show that for every \(j\) , \(\mathrm{Tor}_{j}^{\mathrm{H}}(\mathrm{H}_{0}, \mathrm{M})\) is an \(\mathrm{H}_{0}\)-graded module of finite length when M is of type C. (One may show this by induction on \(j\) . If \(j = 0\) , use a) and b). In the general case, use an exact sequence \(0 \to \mathbf{N} \to \mathbf{H} \otimes_{\mathrm{H}_{0}} \mathbf{M}_{0} \to \mathbf{M} \to 0\) , the induction hypothesis and b).)
\item
  For every H-module M of type C, we set
\end{enumerate}

\[
e (\mathbf {M}) = \sum_ {p \in \mathbf {Z}} \sum_ {j = 0} ^ {\infty} (- 1) ^ {j} [ \operatorname{Tor} _ {j} ^ {\mathrm{H}} (\mathrm{H} _ {0}, \mathbf {M}) _ {p} ] \mathrm{T} ^ {p},
\]

  thus obtaining an element of \(\mathbf{K}(\mathcal{C}_{0})[\mathrm{T},\mathrm{T}^{-1}]\) where \(\mathbf{K}(\mathcal{C}_{0})\) is the Grothendieck group of \(\mathrm{H}_{0}\)-modules of finite length. Show that \(\mathbf{M}\mapsto e(\mathbf{M})\) is additive on exact sequences of graded H-modules, and consequently \(\mathbf{M}\mapsto e(\mathbf{M})\) determines a group homomorphism

\[
e: \mathrm{K} (\mathfrak {C}) \to \mathrm{K} (\mathfrak {C} _ {0}) [ \mathrm{T}, \mathrm{T} ^ {- 1} ].
\]

For every \(H_{0}\) -graded module P of finite length, we set

\[
\tau (\mathrm{P}) = [ \mathrm{H} \otimes_ {\mathrm{H} _ {0}} \mathrm{P} ] \in \mathrm{K} (\mathbb {C}).
\]

Show that \(\mathbf{P} \mapsto \tau(\mathbf{P})\) is additive on exact sequences and that consequently \(\mathbf{P} \mapsto \tau(\mathbf{P})\) determines a homomorphism between the corresponding Grothendieck groups

\[
\tau : \mathrm{K} (\mathbb {C} _ {0}) [ \mathrm{T}, \mathrm{T} ^ {- 1} ] \rightarrow \mathrm{K} (\mathbb {C}).
\]

Show that \(e\) and \(\tau\) are mutually inverse group isomorphisms. (One will first show that \(e \circ \tau\) is the identity. To prove the assertion, it therefore suffices to show that \(\tau\) is surjective, that is, to show that for every H-module M of type C, {[}M{]} belongs to the image of \(\tau\) . One will note that M admits a finite filtration whose successive quotients are annihilated by a maximal ideal of \(H_0\) and that an H-module M of type C, annihilated by a maximal ideal m of \(H_0\) , admits a finite graded resolution by finitely generated free H/mH-modules (A, X, p.~56).)

\begin{enumerate}
\def\labelenumi{\alph{enumi})}
\setcounter{enumi}{4}
\tightlist
\item
  Let M be an H-module of type C. Show that in K(C) one has
\end{enumerate}

\[
[ \mathbf {M} ] = \sum_ {i \geqslant 0} (- 1) ^ {i} [ \mathrm{H} \otimes_ {\mathrm{H} _ {0}} \operatorname{Tor} _ {i} ^ {\mathrm{H}} (\mathrm{H} _ {0}, \mathbf {M}) ].
\]

\begin{enumerate}
\def\labelenumi{\alph{enumi})}
\setcounter{enumi}{5}
\tightlist
\item
  Let \(\Gamma\) be a commutative group and \(\Phi : K(\mathbb{C}) \to \Gamma\) a group homomorphism. For every \(H_0\) -module P of finite length, set
\end{enumerate}

\[
\psi (\mathrm{P}) = \Phi ([ \mathrm{H} \otimes_ {\mathrm{H} _ {0}} \mathrm{P} ]).
\]

Show that for every H-module M of type C, one has

\[
\Phi ([ \mathbf {M} ]) = \sum_ {j \geqslant 0} (- 1) ^ {j} \psi (\operatorname{Tor} _ {j} ^ {\mathrm{H}} (\mathrm{H} _ {0}, \mathbf {M})).
\]

In particular, with the notations of remark 4 of § 4, no. 2, show that one has

\[
\mathrm{Q} _ {\mathbf {M}} = \sum_ {j = 0} ^ {r} (- 1) ^ {j} \mathrm{P} _ {\mathrm{T} _ {j}},
\]

and consequently that

\[
c _ {\mathbf {M}} = \sum_ {j = 0} ^ {r} (- 1) ^ {j} \mathrm{long} _ {\mathrm{H} _ {0}} (\mathrm{Tor} _ {j} ^ {\mathrm{H}} (\mathrm{H} _ {0}, \mathbf {M})).
\]

\begin{enumerate}
\def\labelenumi{\alph{enumi})}
\setcounter{enumi}{6}
\tightlist
\item
  Using the preceding results, give a new proof of th. 1 of § 4, no. 1.
\item
  Let Γ be a commutative group. Describe all additive maps M ↦ Φ(M) on H-modules of type C with values in Γ, in terms of maps from Specmax(H₀) × Z to Γ.
\item
  Assume that H₀ is a field. Let A be a ring. Describe all additive maps \(M\mapsto C(M)\) with values in \(A\) such that \(C(M)\cdot C(N) = \sum_j(-1)^j C(\operatorname{Tor}_{j}^{\mathrm{H}}(M,N))\) for every pair M and N of finitely generated graded H-modules.
\end{enumerate}

\begin{enumerate}
\def\labelenumi{\arabic{enumi})}
\setcounter{enumi}{4}
\tightlist
\item
  \begin{enumerate}
  \def\labelenumii{\alph{enumii})}
  \tightlist
  \item
    Let H be a graded ring of type N, such that \(H_{0}\) is a field and H is a finitely generated \(H_{0}\) -algebra. Let M and N be two finitely generated graded H-modules. Set \(T_{i} = \operatorname{Tor}_{i}^{\mathrm{H}}(\mathbf{M}, \mathbf{N})\) . Prove the equality
  \end{enumerate}
\end{enumerate}

\[
\sum_ {i = 0} ^ {\infty} (- 1) ^ {i} \mathrm{P} _ {\mathrm{T} _ {i}} = \frac {\mathrm{P} _ {\mathrm{M}} \cdot \mathrm{P} _ {\mathrm{N}}}{\mathrm{P} _ {\mathrm{H}}}.
\]

\begin{enumerate}
\def\labelenumi{\alph{enumi})}
\setcounter{enumi}{1}
\tightlist
\item
  Let A be a Noetherian local ring, and M and N two A-modules. Show that if, for \(m \in Z\) and \(i \geqslant 1\) , \(\operatorname{Tor}_{i}^{\operatorname{gr}_{m}(A)}(\operatorname{gr}_{m}(M), \operatorname{gr}_{m}(N))\) is zero, then one has \(\operatorname{Tor}_{i}^{A}(M, N) = 0\) for \(i \geqslant 1\) and \(H_{M \otimes_{A} N, m} = \frac{H_{M, m} \cdot H_{N, m}}{H_{A, m}}\) for \(m \in Z\) .
\end{enumerate}

\begin{enumerate}
\def\labelenumi{\arabic{enumi})}
\setcounter{enumi}{5}
\tightlist
\item
  Let H be a graded ring of type Z, such that \(H_{n} = \{0\}\) for n \textless{} 0, that \(H_{0}\) is an Artinian local ring and that H is a finitely generated \(H_{0}\) -algebra. Let L. be a bounded complex of finitely generated free graded H-modules (cf. A, X, p.~56). Thus for each i, \(L_{i}\) is isomorphic to an H-module of the form \(\bigoplus_{j=1}^{r_{i}} H(-n_{ij})\) , with \(n_{ij} \in Z\) , where \(H(-k)\) is the graded H-module such that \(H(-k)_{n} = H_{n-k}\) . For every finitely generated graded H-module M, denote by \(G_{M} \in Q[T]\) the unique polynomial such that there exists an integer \(n_{1} \geqslant 0\) satisfying
\end{enumerate}

\[
\mathrm{G} _ {\mathrm{M}} (n) = \sum_ {i <   n} \operatorname{long} _ {\mathrm{H} _ {0}} (\mathrm{M} _ {i}), \quad \text {for} \quad n \geqslant n _ {1}.
\]

Prove the relation

\[
\sum_ {i \geqslant 0} (- 1) ^ {i} \mathrm{G} _ {\mathrm{H} _ {i} (\mathrm{L.})} = \sum_ {k \geqslant 0} \frac {(- 1) ^ {k}}{k !} \rho_ {k} \frac {\partial^ {k} \mathrm{G} _ {\mathrm{H}}}{\partial \mathrm{T} ^ {k}}
\]

where \(\rho_{k} = \sum_{i\geqslant 0}(-1)^{i}\sum_{j = 1}^{r_{i}}n_{ij}\) , the \(n_{ij}\) being the degrees associated with \(\mathbf{L}_i\) , and where \(\mathrm{H}_i(\mathrm{L.})\) denotes the homology of the complex \(\mathrm{L.}\).

¶ 7) a) Let l be an integer ≥ 1 and n ∈ N. Show that there exists a unique decreasing sequence of integers: \(a_{l,l}(n) \geqslant a_{l,l-1}(n) \geqslant \cdots \geqslant a_{l,1}(n) \geqslant 0\) such that

\[
n = \sum_ {i = 1} ^ {l} \binom {a _ {l, i} (n) + i - 1} {i}.
\]

Show that one has \(n \leqslant m\) if and only if \((a_{l,l}(n), a_{l,l-1}(n), ..., a_{l,1}(n))\) is less than or equal to \((a_{l,l}(m), ..., a_{l,1}(m))\) for the lexicographic order.

One sets

\[
\partial_ {l} (n) = \sum_ {i = 1} ^ {l} \binom {a _ {l, i} (n) + i} {i + 1},
\]

and, for \(l \geqslant 2\) , denote by \(\partial_l^*(n)\) the smallest integer \(p\) such that \(\partial_{l-1}(p) \geqslant n\) .

Compute \(a_{l+1,i}(\partial_l(n)), a_{l-1,i}(\partial_l^*(n))\) in terms of the \(a_{l,i}(n)\) .

\begin{enumerate}
\def\labelenumi{\alph{enumi})}
\setcounter{enumi}{1}
\tightlist
\item
  One says that a map H:N → N is a Macaulay function if
\end{enumerate}

\[
\left\{ \begin{array}{l l} \mathrm{H} (0) = 1, \\ \mathrm{H} (l + 1) \leqslant \partial_ {l} (\mathrm{H} (l)) & \text {for all} \quad l \geqslant 1. \end{array} \right.
\]

Show that this is equivalent to saying that

\[
\left\{ \begin{array}{l} \mathrm{H} (0) = 1, \\ \partial_ {l} ^ {*} \mathrm{H} (l) \leqslant \mathrm{H} (l - 1) \quad \text {for all} \quad l \geqslant 2. \end{array} \right.
\]

Let H : N → N be a Macaulay function. Show that only two cases are possible: α) There exists r ∈ N such that H(j) = 0 for all j ≥ r.

β) There exists \(r \in \mathbf{N}\) and a finite sequence of integers

\[
b _ {0} \geqslant b _ {1} \geqslant \dots \geqslant b _ {r} \geqslant 0
\]

such that for all \(j \geqslant r\) , one has

\[
(*) \quad \mathrm{H} (j) = \sum_ {n = 0} ^ {r} \binom {b _ {n} - n + j} {b _ {n}}.
\]

Show that \(r\) and the sequence \(b_0 \geqslant b_1 \geqslant \cdots \geqslant b_r \geqslant 0\) are uniquely determined by the function \(H\). We set \(d(H) = 0\) in case \(\alpha\) and \(d(H) = b_0 + 1\) in case \(\beta\), so that \(d(H) - 1\) is the degree of the polynomial (*) in the variable \(j\). Let \(a(H) \frac{j^{d(H)-1}}{(d(H) - 1)!}\) be the term of highest degree of the polynomial (*). Show that \(a(H)\) is an integer \(> 0\).

Show that a polynomial in one variable \(j\) with coefficients in \(\mathbf{Q}\) , which takes integer values on integers and is strictly positive for \(j\) large enough, is not necessarily of the form (*).

\begin{enumerate}
\def\labelenumi{\alph{enumi})}
\setcounter{enumi}{2}
\tightlist
\item
  Let H: N → N be a Macaulay function. We set
\end{enumerate}

\[
\mathrm{S} _ {\mathrm{H}} (\mathrm{T}) = \sum_ {n = 0} ^ {\infty} \mathrm{H} (n) \mathrm{T} ^ {n},
\]

so that \(S_{H}(T) \in Z[[T]]\) . Show that \(S_{H}(T) \in Z[T]\) or else that there exists a finite sequence of integers

\[
c _ {0} \geqslant c _ {1} \geqslant \dots \geqslant c _ {r} > 0
\]

uniquely determined by H, and a polynomial \(P \in Z[T]\) such that one has

\[
(* *) \mathrm{S} _ {\mathrm{H}} (\mathrm{T}) = \sum_ {n = 0} ^ {r} \frac {\mathrm{T} ^ {n}}{(1 - \mathrm{T}) ^ {c _ {n}}} + \mathrm{P}.
\]

Show that \(c_0 = d(\mathbf{H})\) and that \(a(\mathbf{H})\) is the number of integers \(c_i\) such that \(c_i = d(\mathbf{H})\) .

The rational fraction \(\sum_{n=0}^{r}\frac{T^n}{(1-T)^{c_n}}\) is called the asymptotic expansion of \(H\). Let \(A_1, ..., A_d\) be integers, with \(A_d \neq 0\), such that the rational fraction

\[
\sum_ {n = 1} ^ {d} \frac {\mathrm{A} _ {n}}{(1 - \mathrm{T}) ^ {n}} - \mathrm{S} _ {\mathrm{H}} (\mathrm{T})
\]

is a polynomial. Show that the sequence \((\mathrm{A}_{n})_{1\leq n\leq d}\) is uniquely determined by, and uniquely determines, the asymptotic expansion of \(H\). In particular, we have \(d=d(\mathrm{H})\), \(\mathrm{A}_{d}=a(\mathrm{H})\). Examine the case \(d(\mathrm{H})=2\) in more detail.

\begin{enumerate}
\def\labelenumi{\alph{enumi})}
\setcounter{enumi}{3}
\tightlist
\item
  Let \(\mathbf{H}:\mathbf{N}\to\mathbf{N}\) be a Macaulay function. Show that the following conditions are equivalent:
\end{enumerate}

\begin{enumerate}
\def\labelenumi{(\roman{enumi})}
\item
  \(\partial_l^*\mathrm{H}(l) = \mathrm{H}(l - 1)\) for every \(l\geqslant 2\)
\item
  H is the smallest Macaulay function having the same asymptotic expansion as H.
\end{enumerate}

Such Macaulay functions are said to be extremal. Let \(c_{0} \geqslant c_{1} \geqslant \cdots \geqslant c_{r} > 0\) be a sequence of integers. Show that there exists one and only one extremal Macaulay function whose asymptotic expansion is \(\sum_{n=0}^{r} \frac{\mathrm{T}^{n}}{(1 - \mathrm{T})^{c_{n}}}\) . Show that, if H is an extremal Macaulay function, then one has \(\mathrm{H}(1) = d(\mathrm{H})\) or \(\mathrm{H}(1) = d(\mathrm{H}) + 1\). Deduce from this that for every Macaulay function H, one has \(\mathrm{H}(1) \geqslant d(\mathrm{H})\) .

\begin{enumerate}
\def\labelenumi{\alph{enumi})}
\setcounter{enumi}{4}
\tightlist
\item
  Let H: N → N be a Macaulay function. Show that the following properties are equivalent:
\end{enumerate}

\begin{enumerate}
\def\labelenumi{(\roman{enumi})}
\tightlist
\item
  we have \(\mathrm{S_H(T)} = \frac{1}{(1 - T)^{d(H)}}\)
\end{enumerate}

\begin{enumerate}
\def\labelenumi{(\roman{enumi})}
\setcounter{enumi}{1}
\item
  \(\mathbf{H}(1) = d(\mathbf{H})\)
\item
  \(\partial_{l}\mathbf{H}(l) = \mathbf{H}(l + 1)\) for every \(l\geqslant 1\)
\end{enumerate}

In particular, any function such that \(\mathrm{H}(1)=d(\mathrm{H})\) is extremal.

\begin{enumerate}
\def\labelenumi{\arabic{enumi})}
\setcounter{enumi}{7}
\tightlist
\item
  Let E be a set. We set \(\mathcal{M}(\mathrm{E}) = \mathbf{N}^{(\mathrm{E})}\) and we denote multiplicatively the composition in \(\mathcal{M}(\mathrm{E})\) . An ideal of \(\mathcal{M}(\mathrm{E})\) is a subset \(I\) of \(\mathcal{M}(\mathrm{E})\) such that \(I\cdot\mathcal{M}(\mathrm{E}) \subset I\) . A staircase of \(\mathcal{M}(\mathrm{E})\) is the complement of an ideal. We denote by \(d: \mathcal{M}(\mathrm{E}) \to \mathbf{N}\) the unique monoid homomorphism such that \(d(x) = 1\) , for all \(x \in \mathrm{E}\) . Let \(\mathrm{A} \subset \mathcal{M}(\mathrm{E})\) . For every \(l \in \mathbf{N}\) , denote by \(\mathrm{A}_l\) the set \(\mathrm{A} \cap d^{-1}(l)\) . We also denote \(\delta \mathrm{A}\) the set of \(m \in \mathcal{M}(\mathrm{E})\) such that, for every element \(x\) of E which divides \(m\) , one has \(m / x \in \mathrm{A}\) and \(d\mathrm{A}\) the set of \(m\) such that there exists an \(x \in \mathrm{E}\) with \(mx \in \mathrm{A}\) .
\end{enumerate}

\begin{enumerate}
\def\labelenumi{\alph{enumi})}
\tightlist
\item
  Show that for every subset A of \(\mathcal{M}(E)\) , one has \(d\delta A \subset A\) and \(A \subset \delta dA\) ; if A and B are two subsets of \(\mathcal{M}(E)\) , the relations \(dA \subset B\) and \(A \subset \delta B\) are equivalent. Show that the following three properties of a subset \(\Delta\) of \(\mathcal{M}(E)\) are equivalent:
\end{enumerate}

\begin{enumerate}
\def\labelenumi{(\roman{enumi})}
\item
  \(\Delta\) is a staircase;
\item
  \(\Delta \supset d\Delta\) ;
\item
  \(\delta \Delta \supset \Delta\) .
\end{enumerate}

\begin{enumerate}
\def\labelenumi{\alph{enumi})}
\setcounter{enumi}{1}
\tightlist
\item
  Let F be a subset of E and identify \(\mathcal{M}(\mathbf{F})\) with its canonical image in \(\mathcal{M}(\mathbf{E})\) . Show that a subset of \(\mathcal{M}(\mathbf{F})\) is a staircase of \(\mathcal{M}(\mathbf{F})\) if and only if it is a staircase of \(\mathcal{M}(\mathbf{E})\) . Show that, for a staircase \(\Delta\) of \(\mathcal{M}(\mathbf{E})\) the following conditions are equivalent:
\end{enumerate}

\begin{enumerate}
\def\labelenumi{(\roman{enumi})}
\item
  \(\Delta_l\) is finite for every \(l \in \mathbf{N}\) ;
\item
  \(\Delta_{1}\) is finite;
\item
  there exists a finite subset F of E such that \(\Delta\subset\mathcal{M}(F)\) .
\end{enumerate}

Such a staircase is said to be of finite type.

\begin{enumerate}
\def\labelenumi{\alph{enumi})}
\setcounter{enumi}{2}
\tightlist
\item
  From now on, we assume that E = N and we set \(\mathcal{M}(E) = \mathcal{M}\) . Let k be a field. We denote by R the graded algebra \(\mathrm{K}^{(\mathcal{M})} = \mathrm{K}[(\mathrm{X}_i)_{i \in \mathbb{N}}]\) . We identify M with the set of monomials in \(\mathrm{K}[(\mathrm{X}_i)_{i \in \mathbb{N}}]\) and, for every subset A of M, we denote by \(k(A)\) the vector subspace of R generated by A. Let I be an ideal of M and \(\Delta\) the complementary staircase. Then \(k(I)\) is an ideal of R and \(k(\Delta)\) is a supplementary subspace of \(k(I)\) . Let \(J \subset R\) be a graded ideal of \(R\). Show that there exists a staircase \(\Delta \subset M\) such that \(k(\Delta)\) is a supplement of \(J\). (Put on \(M\) the reverse lexicographic order for which one has \(X_1^{a_1} \ldots X_n^{a_n} \prec X_1^{b_1} \ldots X_n^{b_n}\) if there exists \(0 \leqslant p \leqslant n\) such that \(a_j = b_j\) for \(j \textgreater{} p\), and \(a_p < b_p\) . Define by induction \(m_1, \ldots, m_p, \ldots\) in \(M\) as follows: for every \(p\), \(m_p\) is the smallest monomial of \(M\) linearly independent of \(m_1, \ldots, m_{p-1}\) modulo \(J\). Let \(\Delta\) be the set of \(m_i\) . Show that \(\Delta\) is a staircase.) Show that such a staircase is of finite type if and only if the algebra \(R/J\) is of finite type over \(k\).
\end{enumerate}

  We use the notation from the previous exercise. We put the lexicographic order on \(M\). Let \(l \in \mathbf{N}\). A subset \(A\) of \(\mathcal{M}_{l}\) is said to be initial if for every \(m \in \mathrm{A}\) , the set

\[
[ m ] = \{m ^ {\prime} \in \mathcal {M} _ {d (m)} | m ^ {\prime} <   m \}
\]

is contained in A.

\begin{enumerate}
\def\labelenumi{\alph{enumi})}
\item
  Let \(l \in \mathbf{N}\) (resp. \(l \in \mathbf{N} - \{0\}\) ) and let \(\mathbf{A} \subset \mathcal{M}_l\) be an initial segment. Show that \(\delta \mathbf{A}\) (resp. \(d\mathbf{A}\) ) is an initial segment of \(\mathcal{M}_{l+1}\) (resp. \(\mathcal{M}_{l-1}\) ). Suppose that \(\mathbf{A}\) has a greatest element \(m = X_1^{a_1} \ldots X_n^{a_n}\) . Determine the greatest element of \(\delta \mathbf{A}\) (resp. \(d\mathbf{A}\) ).
\item
  Let \(l \in \mathbf{N} - \{0\}\) and let \(\mathbf{A} \subset \mathcal{M}_l\) be a nonempty finite initial segment. Let \(\mathrm{X}_{a_l}\) be the largest variable such that \([\mathrm{X}_{a_l}^l] \subset \mathbf{A}\) . Show that every \(m\) in \(\mathbf{A} - [\mathrm{X}_{a_l}^l]\) is equal to 1 or divisible by \(\mathrm{X}_{a_l + 1}\) and that \(\mathrm{A}' = \{m \in \mathcal{M}_{l-1}|m\mathrm{X}_{a_l + 1} \in \mathbf{A} - [\mathrm{X}_{a_l}^l]\}\) is an initial segment of \(\mathcal{M}_{l-1}\). Deduce from the preceding an expression for \(\operatorname{Card}([X_1^{b_1} ... X_n^{b_n}])\) .
\item
  Deduce from the preceding that, for every \(l \in \mathbf{N} - \{0\}\) and every finite initial segment A of \(\mathcal{M}_l\) , one has \(\text{Card}(\delta\text{A}) = \partial_l(\text{Card}(\text{A}))\) , where \(\partial_l\) was defined in exerc. 7. Compute analogously \(\text{Card}(d\text{A})\) .
\item
  For every \(l \in \mathbf{N}\) and every \(n \in \mathbf{N}\) , set
\end{enumerate}

\[
\mu (n) = \inf_{\substack{\mathbf{A}\subset \mathcal{M}_{l}\\ \operatorname{Card}(\mathbf{A}) = n}}\operatorname{Card}(d\mathbf{A}),
\]

and denote by \(m_{n}\) the \(n\) -th monomial of \(\mathcal{M}_{l}\) for the lexicographic order (we set \(m_{0}=1\) ).

For every finite subset A of \(\mathcal{M}_{l}\) , denote by CA the initial segment of \(\mathcal{M}_{l}\) such that \(\mathrm{Card(A)} = \mathrm{Card(CA)}\) . Show that the following three assertions are equivalent:

(MC1) for every \(l \in \mathbf{N}\) and every \(\mathbf{A} \subset \mathcal{M}_l\) , one has \(\mathbf{C}\delta \mathbf{A} \subset \delta \mathbf{CA}\) ;

(MC2) for every \(l \in \mathbf{N}\) and every \(\mathbf{A} \subset \mathcal{M}_l\) , one has \(d\mathbf{CA} \subset \mathbf{CdA}\) ;

(MC3) for every \(l \in \mathbf{N}\) and every \(n \in \mathbf{N}\) , one has \(\operatorname{Card}(d[m_n]) = \mu(n)\) .

\begin{enumerate}
\def\labelenumi{\arabic{enumi})}
\setcounter{enumi}{9}
\tightlist
\item
  We use the notation of the two preceding exercises. Let \(l \in \mathbf{N}\) and \(\mathbf{A} \subset \mathcal{M}_l\) . For every pair \((i, d)\) one sets
\end{enumerate}

\(\mathrm{A}_{(i;d)} = \{m \in \mathrm{A} | m \text{ is divisible by } \mathrm{X}_i^d \text{ and is not divisible by } \mathrm{X}_i^{d+1}\}\)

and we denote by \(C_{(i;d)}A\) the set of \(\operatorname{Card}(A_{(i;d)})\) first elements of \((\mathcal{M}_{l})_{(i;d)}\) (for the lexicographic order). We set \(C_{i}A = \bigcup C_{(i;d)}A\) .

\begin{enumerate}
\def\labelenumi{\alph{enumi})}
\tightlist
\item
  For every \(m \in \mathcal{M}_{l}\) , denote by \(\pi(m)\) the predecessor in \(\mathcal{M}_{l}\) of \(m\) for the reverse lexicographic order, when it exists (i.e., \(m \neq X_{1}^{l}\) ). Let \(i\) be an integer such that \(1 < i\), and let \(a_{i}\) be a strictly positive integer. Show that we have
\end{enumerate}

\[
\begin{array}{c} \pi (X _ {i} ^ {a _ {i}} \dots X _ {p} ^ {a _ {p}}) = X _ {i - 1} X _ {i} ^ {a _ {i} - 1} \dots X _ {p} ^ {a _ {p}} \\ \pi (X _ {1} ^ {a _ {1}} X _ {i} ^ {a _ {i}} \dots X _ {p} ^ {a _ {p}}) = X _ {i - 1} ^ {a _ {i - 1} + 1} X _ {i} ^ {a _ {i} - 1} \dots X _ {p} ^ {a _ {p}}. \end{array}
\]

\begin{enumerate}
\def\labelenumi{\alph{enumi})}
\setcounter{enumi}{1}
\item
  Let \(\mathbf{A} \subset \mathcal{M}_l\) and \(p \geqslant 4\) such that we have \(\mathbf{A} \subset [\mathbf{X}_p^l]\) (i.e., such that the monomials belonging to \(\mathbf{A}\) involve only the variables \(\mathbf{X}_j\) , \(1 \leqslant j \leqslant p\) ). We assume that for every integer \(i\) such that \(1 \leqslant i \leqslant p\) , one has \(\mathbf{A} = \mathbf{C}_i\mathbf{A}\) . Show that \(\mathbf{A} = \mathbf{CA}\) (notations of exerc. 9, d)).
\item
  Let \(\mathbf{A} \subset \mathcal{M}_l\) such that \(\mathbf{A} \subset [\mathbf{X}_3^l]\) and that \(\mathbf{A} = \mathbf{C}_i\mathbf{A}\) for \(i = 1, 2, 3\) . Put \(\mathbf{A} = \bigcup_{k=0} \mathbf{B}_k\mathbf{X}_3^{k}\) with \(\mathbf{B}_k \subset [\mathbf{X}_2^{l-k}]\) and \(b(k) = \text{Card}(\mathbf{B}_k)\) . Show that \(\mathbf{B}_k\) is an initial segment of \(\mathcal{M}_{l-k}\) and that \(0 \leqslant b(0) \leqslant l + 1\) , \(b(k+1) < b(k)\) if \(b(k) \neq 0\) , and that \(k \mapsto b(k)\) is decreasing. Examine the converse of this assertion. Let \(\mu\) be the number of integers \(k\) such that \(b(k) = l + 1 - k\) . Show that \(\text{Card}(d\mathbf{A}) = (\text{Card}(\mathbf{A})) - \mu\) .
\end{enumerate}

¶ 11) We propose to prove assertions MC1, MC2, MC3 of exercise 9, d) (Macaulay's theorem). Let \(l \in \mathbf{N}\) , \(\mathbf{A} \subset \mathcal{M}_l\) , and let \(p\) be the smallest integer such that \(\mathbf{A} \subset [\mathbf{X}_p^l]\) . We shall prove MC2 by induction on \(p\) .

\begin{enumerate}
\def\labelenumi{\alph{enumi})}
\item
  Prove the theorem for \(p = 1\) and \(p = 2\) .
\item
  Using the induction hypothesis, show that for all \(i \leqslant p\) , \(dC_{i}A \subset C_{i}dA\) (notations of Exerc. 8).
\item
  One sets \(\mathbf{A}^1 = \mathbf{A}\) , \(\mathbf{A}^{j+1} = \mathbf{C}_i\mathbf{A}^j\) where \(i \equiv j \mod p\) and \(1 \leqslant i \leqslant p\) . Show that for \(j\) sufficiently large one has \(\mathbf{A}^{j+1} = \mathbf{A}^j\) . (For every \(a \in \mathcal{M}_l\) , let \(n(a) \in \mathbf{N}\) denote the rank of \(a\) in the lexicographic order and set \(n(\mathbf{A}) = \sum_{a \in \mathbf{A}} n(a)\) . Show that \(n(\mathbf{A}^{j+1}) \leqslant n(\mathbf{A}^j)\) and that \(n(\mathbf{A}^{j+1}) = n(\mathbf{A}^j)\) if and only if \(\mathbf{A}^{j+1} = \mathbf{A}^j\) .)
\item
  Prove the theorem in the general case. (A subset \(A \subset [X_{p}^{l}] \subset M_{l}\) is said to be minimal if \(\operatorname{Card}(dA) \leqslant \operatorname{Card}(dA')\) for the subsets \(A' \subset [X_{p}^{l}]\) such that \(\operatorname{Card}(A') = \operatorname{Card}(A)\). It is a matter of showing that if A is minimal, then CA is minimal. Let A be minimal. Using b), show that \(C_{i}A\) is minimal for all \(i \leqslant p\) . Deduce, using the notation of c), that \(A^{j}\) is minimal for all j. Deduce from c) that there exists a minimal B such that \(\operatorname{Card}(B) = \operatorname{Card}(A)\) and \(C_{i}B = B\) for \(1 \leqslant i \leqslant p\) . In the case p = 3, use exerc. 10, c) to conclude. In the case \(p \geqslant 4\) using Exerc. 10, b) to deduce that \(B = CA\) .
\end{enumerate}

¶ 12) Let H : N → N be a map and let k be a field. We propose to prove the equivalence of the following properties:

\begin{enumerate}
\def\labelenumi{(\roman{enumi})}
\item
  H is a Macaulay function (p.~90, exerc. 7).
\item
  There exists a staircase of finite type (p.~92, exerc. 8) \(\Delta \subset \mathcal{M}\) such that for all \(l \in \mathbf{N}\) , one has \(\mathrm{H}(l) = \mathrm{Card}(\Delta_l)\) .
\item
  There exists a graded \(k\)-algebra \(\mathbf{A}\) , of type \(\mathbf{N}\) generated by a finite number of its elements of degree 1, such that for every \(l\) , one has \(\mathrm{H}(l) = \dim_k A_l\) .
\end{enumerate}

\begin{enumerate}
\def\labelenumi{\alph{enumi})}
\item
  To prove (ii) \(\Leftrightarrow\) (iii), we will use exerc. 8, c).
\item
  We prove (i) \(\Rightarrow\) (ii). Let H be a Macaulay function. For every \(l\) , let \(\Delta_l \subset \mathcal{M}_l\) be the initial part such that \(\operatorname{Card}(\Delta_l) = \mathrm{H}(l)\) and let us set \(\Delta = \bigcup_{l} \Delta_l\) . To show that \(\Delta\) is a staircase, we will use exerc. 9, c).
\item
  We prove (ii) \(\Rightarrow\) (i). Let \(\Delta\) be a staircase. To show that \(\mathbf{H}: l \mapsto \operatorname{Card}(\Delta_l)\) is a Macaulay function, we will use Macaulay's theorem (MC1) (exerc. 11) and exerc. 9, c) \(^{1}\) .
\end{enumerate}
% semantic-fragments: legacy-input-seam
\relax{}
